% Book-specific typography builds on the reviewed course template.
% Shared lecture-note preamble for Big Data and Management Decision-Making.
% Chapter files follow latex_config.md and remain independently compilable.

\usepackage[a4paper,margin=2.6cm]{geometry}
\usepackage{amsmath,amssymb,amsthm,mathtools,bm}
\usepackage{xcolor}
\usepackage[most]{tcolorbox}
\usepackage{booktabs,longtable,array}
\usepackage{enumitem}
\usepackage{xurl}
\usepackage{hyperref}
\usepackage{listings}
\usepackage{caption}
\usepackage{tikz}
\usepackage{graphicx,pgfplots}
\pgfplotsset{compat=1.18}
\usetikzlibrary{arrows.meta,positioning,shapes.geometric,calc,fit}

\setCJKmainfont[
  UprightFont=FandolSong-Regular.otf,
  BoldFont=FandolSong-Regular.otf,
  ItalicFont=FandolKai-Regular.otf,
  BoldItalicFont=FandolKai-Regular.otf
]{FandolSong-Regular.otf}
\setCJKsansfont[
  UprightFont=FandolSong-Regular.otf,
  BoldFont=FandolSong-Regular.otf
]{FandolSong-Regular.otf}
\setCJKmonofont[
  UprightFont=FandolSong-Regular.otf,
  BoldFont=FandolSong-Regular.otf
]{FandolSong-Regular.otf}
\IfFontExistsTF{Menlo}{\setmonofont{Menlo}}{\setmonofont{lmmono10-regular.otf}}
\linespread{1.13}
\setlength{\emergencystretch}{2em}
\raggedbottom
\setcounter{tocdepth}{1}
\captionsetup{font=small,labelfont=normalfont}
\tikzset{
  bdmnode/.style={draw=BDMTeal,fill=BDMLight,rounded corners=2pt,
    align=center,text width=2.65cm,minimum height=1cm,font=\small},
  bdmarrow/.style={-{Latex[length=2mm]},thick,draw=BDMGray}
}

\hypersetup{
  colorlinks=true,
  linkcolor=blue!60!black,
  citecolor=blue!60!black,
  urlcolor=blue!60!black,
  bookmarksdepth=2,
  pdfauthor={大数据与管理决策基础},
  pdfsubject={大数据与管理决策基础课程讲义}
}

\ctexset{
  chapter = {
    format = \huge\mdseries,
    name = {第,讲},
    number = \arabic{chapter},
    aftername = \quad
  },
  section = {format = \Large\mdseries},
  subsection = {format = \large\mdseries},
  subsubsection = {format = \normalsize\mdseries}
}

\setlist[itemize]{leftmargin=2em,itemsep=0.25em,topsep=0.25em}
\setlist[enumerate]{leftmargin=2em,itemsep=0.25em,topsep=0.25em}

\definecolor{BDMBlue}{RGB}{22,44,73}
\definecolor{BDMTeal}{RGB}{22,119,119}
\definecolor{BDMGray}{RGB}{76,84,95}
\definecolor{BDMLight}{RGB}{245,247,250}
\definecolor{BDMAmber}{RGB}{181,111,35}
\definecolor{CodeBg}{RGB}{248,248,248}

\lstdefinestyle{pythonstyle}{
  basicstyle=\ttfamily\small,
  backgroundcolor=\color{CodeBg},
  keywordstyle=\color{BDMBlue},
  commentstyle=\color{BDMGray},
  stringstyle=\color{BDMAmber},
  showstringspaces=false,
  breaklines=true,
  frame=single,
  rulecolor=\color{black!20},
  columns=fullflexible
}

\lstdefinestyle{sqlstyle}{
  language=SQL,
  basicstyle=\ttfamily\small,
  backgroundcolor=\color{CodeBg},
  keywordstyle=\color{BDMBlue},
  commentstyle=\color{BDMGray},
  stringstyle=\color{BDMAmber},
  showstringspaces=false,
  breaklines=true,
  frame=single,
  rulecolor=\color{black!20},
  columns=fullflexible
}

\lstset{style=pythonstyle}
\renewcommand{\lstlistingname}{代码}
\renewcommand{\bibname}{参考文献}

\theoremstyle{definition}
\newtheorem{definition}{定义}[chapter]
\newtheorem{example}[definition]{例}
\newtheorem{exercise}[definition]{习题}
\newtheorem{convention}[definition]{约定}

\theoremstyle{remark}
\newtheorem{remark}[definition]{评注}

\theoremstyle{plain}
\newtheorem{theorem}[definition]{定理}
\newtheorem{proposition}[definition]{命题}
\newtheorem{lemma}[definition]{引理}
\newtheorem{corollary}[definition]{推论}

\tcolorboxenvironment{definition}{
  colback=blue!2!white,
  colframe=blue!45!black,
  boxrule=0.4pt,
  arc=1.2mm,
  breakable
}

\tcolorboxenvironment{theorem}{
  colback=orange!3!white,
  colframe=orange!55!black,
  boxrule=0.4pt,
  arc=1.2mm,
  breakable
}

\tcolorboxenvironment{proposition}{
  colback=orange!2!white,
  colframe=orange!50!black,
  boxrule=0.4pt,
  arc=1.2mm,
  breakable
}

\tcolorboxenvironment{lemma}{
  colback=orange!2!white,
  colframe=orange!45!black,
  boxrule=0.4pt,
  arc=1.2mm,
  breakable
}

\tcolorboxenvironment{corollary}{
  colback=orange!2!white,
  colframe=orange!45!black,
  boxrule=0.4pt,
  arc=1.2mm,
  breakable
}

\newtcolorbox{chapterbox}{
  colback=gray!5!white,
  colframe=gray!55!black,
  boxrule=0.4pt,
  arc=1.2mm,
  breakable
}

\newtcolorbox{modernbox}[1][应用接口]{
  colback=green!3!white,
  colframe=green!45!black,
  boxrule=0.4pt,
  arc=1.2mm,
  title={#1},
  fonttitle=\mdseries,
  breakable
}

\newcommand{\R}{\mathbb R}
\newcommand{\N}{\mathbb N}
\newcommand{\E}{\mathbb E}
\newcommand{\Prob}{\mathbb P}
\newcommand{\dd}{\,\mathrm d}
\newcommand{\eps}{\varepsilon}
\newcommand{\abs}[1]{\left\lvert #1\right\rvert}
\newcommand{\norm}[1]{\left\lVert #1\right\rVert}
\newcommand{\argmin}{\operatorname*{arg\,min}}
\newcommand{\argmax}{\operatorname*{arg\,max}}

% Editorial overrides are local to the book; historical lectures stay unchanged.
\usepackage{colortbl,needspace,microtype}
\definecolor{BDMBlue}{HTML}{263A3D}
\definecolor{BDMTeal}{HTML}{216B68}
\definecolor{BDMGray}{HTML}{505A5C}
\definecolor{BDMAmber}{HTML}{A54932}
\definecolor{BDMLight}{HTML}{F4F7F6}
\definecolor{BDMRule}{HTML}{A9B8B6}
\definecolor{BDMTableHead}{HTML}{EAF0EF}
\definecolor{CodeBg}{HTML}{F5F7F7}

\setmonofont{IBMPlexMono-Regular.otf}[
  BoldFont=IBMPlexMono-Medium.otf,
  ItalicFont=IBMPlexMono-Italic.otf,
  BoldItalicFont=IBMPlexMono-MediumItalic.otf,
  Scale=0.88,Ligatures=NoCommon]
\setsansfont{texgyreheros-regular.otf}[
  BoldFont=texgyreheros-bold.otf,ItalicFont=texgyreheros-italic.otf,
  BoldItalicFont=texgyreheros-bolditalic.otf,Scale=0.92]
\setCJKsansfont[
  UprightFont=FandolHei-Regular.otf,BoldFont=FandolHei-Regular.otf
]{FandolHei-Regular.otf}

\ctexset{
  part={format=\Huge\mdseries\color{BDMBlue}},
  chapter={format=\huge\mdseries\color{BDMBlue}},
  section={format=\Large\mdseries\color{BDMBlue}},
  subsection={format=\large\mdseries},
  subsubsection={format=\normalsize\mdseries}
}
\setlength{\parskip}{0.16em plus 0.06em minus 0.03em}
\setlength{\textfloatsep}{16pt plus 3pt minus 3pt}
\setlength{\intextsep}{14pt plus 3pt minus 3pt}
\setlength{\floatsep}{12pt plus 2pt minus 2pt}
\renewcommand{\topfraction}{0.85}
\renewcommand{\bottomfraction}{0.75}
\renewcommand{\textfraction}{0.12}
\renewcommand{\floatpagefraction}{0.8}
\setcounter{topnumber}{3}
\setcounter{totalnumber}{4}
\makeatletter
\setlength{\@fptop}{16pt}
\setlength{\@fpsep}{22pt}
\setlength{\@fpbot}{0pt plus 1fil}
\makeatother
\clubpenalty=10000
\widowpenalty=10000
\displaywidowpenalty=10000
\captionsetup{font={small},labelfont={normalfont,color=BDMTeal},
  textfont={normalfont,color=BDMGray},labelsep=quad,skip=8pt,
  justification=raggedright,singlelinecheck=true}
\captionsetup[table]{position=top}
\captionsetup[lstlisting]{position=top}

\newcolumntype{P}[1]{>{\raggedright\arraybackslash}p{#1}}
\newcommand{\booktableformat}{%
  \small\renewcommand{\arraystretch}{1.22}%
  \setlength{\tabcolsep}{5pt}\arrayrulecolor{BDMRule}%
  \setlength{\heavyrulewidth}{0.65pt}%
  \setlength{\lightrulewidth}{0.3pt}}
\AtBeginEnvironment{table}{\booktableformat}
\AtBeginEnvironment{longtable}{\booktableformat}
\setlength{\LTpre}{12pt}
\setlength{\LTpost}{12pt}

\lstdefinestyle{bookcode}{
  basicstyle=\ttfamily\small\color{BDMBlue},
  backgroundcolor=\color{CodeBg},
  keywordstyle=\color{BDMTeal},
  commentstyle=\color{BDMGray},
  stringstyle=\color{BDMAmber},
  showstringspaces=false,showspaces=false,showtabs=false,
  keepspaces=true,columns=fullflexible,tabsize=4,
  breaklines=true,breakatwhitespace=false,
  frame=l,framerule=0.7pt,rulecolor=\color{BDMRule},
  xleftmargin=7pt,xrightmargin=4pt,framexleftmargin=5pt,
  framesep=5pt,aboveskip=13pt,belowskip=13pt,
  captionpos=t,numbers=none,upquote=true}
\lstdefinestyle{pythonstyle}{style=bookcode,language=Python}
\lstdefinestyle{sqlstyle}{style=bookcode,language=SQL}
\lstset{style=pythonstyle}
\BeforeBeginEnvironment{lstlisting}{\Needspace{9\baselineskip}}

\tcbset{bookstatement/.style={enhanced,breakable,
  colback=BDMLight,coltext=black,frame hidden,boxrule=0pt,arc=0pt,
  borderline west={0.9pt}{0pt}{BDMTeal},
  left=9pt,right=9pt,top=6pt,bottom=6pt,
  before skip=11pt,after skip=11pt}}
\RemoveFromHook{env/definition/before}[tcolorbox]
\RemoveFromHook{env/definition/after}[tcolorbox]
\RemoveFromHook{env/theorem/before}[tcolorbox]
\RemoveFromHook{env/theorem/after}[tcolorbox]
\RemoveFromHook{env/proposition/before}[tcolorbox]
\RemoveFromHook{env/proposition/after}[tcolorbox]
\RemoveFromHook{env/lemma/before}[tcolorbox]
\RemoveFromHook{env/lemma/after}[tcolorbox]
\RemoveFromHook{env/corollary/before}[tcolorbox]
\RemoveFromHook{env/corollary/after}[tcolorbox]
\tcolorboxenvironment{definition}{bookstatement}
\tcolorboxenvironment{theorem}{bookstatement,borderline west={0.9pt}{0pt}{BDMAmber}}
\tcolorboxenvironment{proposition}{bookstatement,borderline west={0.9pt}{0pt}{BDMAmber}}
\tcolorboxenvironment{lemma}{bookstatement}
\tcolorboxenvironment{corollary}{bookstatement}
\renewtcolorbox{chapterbox}{bookstatement}
\renewtcolorbox{modernbox}[1][应用研究]{bookstatement,
  title={#1},coltitle=BDMTeal,fonttitle=\mdseries,
  attach title to upper={\par\smallskip},colbacktitle=BDMLight}

\tikzset{
  bdmnode/.style={draw=BDMRule,fill=BDMLight,line width=0.55pt,
    rounded corners=1pt,align=center,text width=2.65cm,
    minimum height=1.02cm,inner sep=5pt,font=\small\sffamily,
    text=BDMBlue},
  bdmarrow/.style={-{Latex[length=1.8mm,width=1.3mm]},
    line width=0.65pt,draw=BDMTeal,shorten >=1pt,shorten <=1pt},
  bdmedge/.style={line width=0.65pt,draw=BDMTeal},
  bdmlabel/.style={font=\footnotesize,fill=white,inner sep=2pt,text=BDMGray},
  bdmvar/.style={circle,draw=BDMRule,fill=BDMLight,
    minimum size=8mm,inner sep=1pt,text=BDMBlue}
}
\pgfplotsset{every axis/.append style={
  axis line style={draw=BDMRule,line width=0.4pt},
  tick style={draw=BDMRule},tick label style={font=\footnotesize},
  label style={font=\small},
  grid=major,major grid style={draw=BDMRule!35,line width=0.25pt},
  legend style={draw=none,fill=white,font=\footnotesize,
    cells={anchor=west}},
  every axis plot/.append style={line width=0.85pt},
  mark size=1.7pt}}

\usepackage{fancyhdr,multicol,makeidx}
\usepackage[backend=biber,style=numeric-comp,sorting=nyt,
  giveninits=true,maxbibnames=8,maxcitenames=2,
  doi=true,isbn=false,url=true,eprint=false,urldate=iso,seconds=true,
  backref=true]{biblatex}
\addbibresource{references.bib}
\makeindex
\renewcommand{\indexname}{主题索引}
\ctexset{
  chapter={name={第,章},beforeskip=18pt,afterskip=24pt},
  section={beforeskip=20pt,afterskip=10pt},
  subsection={beforeskip=14pt,afterskip=7pt}
}
\setcounter{tocdepth}{1}
\setcounter{secnumdepth}{2}
\setlength{\headheight}{16pt}
\setlength{\headsep}{18pt}
\setlength{\footskip}{26pt}
\pagestyle{fancy}
\fancyhf{}
\fancyhead[L]{\footnotesize\color{BDMGray} 大数据与管理决策基础}
\fancyhead[R]{\footnotesize\color{BDMGray}\nouppercase{\leftmark}}
\fancyfoot[C]{\small\color{BDMGray}\thepage}
\renewcommand{\headrulewidth}{0.2pt}
\renewcommand{\chaptermark}[1]{\markboth{\thechapter\quad #1}{}}
\fancypagestyle{plain}{
  \fancyhf{}\fancyfoot[C]{\thepage}\renewcommand{\headrulewidth}{0pt}
}
\hypersetup{
  pdftitle={大数据与管理决策基础：数据分析、统计推断与管理决策},
  pdfauthor={孙振宇；东北大学数学与统计学院},
  pdfsubject={十章统一教材；公司治理、统计学习、因果推断与决策系统},
  linkcolor=BDMBlue,citecolor=BDMBlue,urlcolor=BDMTeal,
  bookmarksnumbered=true,bookmarksopen=false
}
\setlength{\bibitemsep}{0.65\baselineskip}
\setcounter{biburllcpenalty}{7000}
\setcounter{biburlucpenalty}{8000}
\setcounter{biburlnumpenalty}{9000}
\renewcommand*{\bibfont}{\small}
\AtBeginBibliography{\interlinepenalty=10000\raggedright}
\DeclareFieldFormat{urldate}{\mkbibbrackets{访问日期：#1}}
\DefineBibliographyStrings{english}{bibliography={参考文献},references={参考文献},
  backrefpage={引用页},backrefpages={引用页}}
\newcommand{\subjectindex}[2]{\index{#1@#2|hyperpage}}
\newcommand{\indexletter}[1]{\par\Needspace{4\baselineskip}\bigskip\noindent\large\mdseries #1\par\nobreak\smallskip\nobreak\normalsize}
\makeatletter
\renewenvironment{theindex}{
  \clearpage\chapter*{主题索引}\markboth{主题索引}{主题索引}
  \addcontentsline{toc}{chapter}{主题索引}
  本索引按中文词目的汉语拼音、英文词目的字母顺序混合排列。
  页码指向概念的主要定义、推导或应用。
  英文缩略语与中文术语分别设入口。\par\medskip
  \begin{multicols}{2}\small\raggedright
  \setlength{\parindent}{0pt}\setlength{\parskip}{3pt}
  \let\item\@idxitem
}{\end{multicols}}
\makeatother
\newcommand{\chapterreading}[1]{%
  \subsection*{文献导读}#1\par
}
